% !TEX root = ../main.tex

\chapter{Conclusion}\label{chapter:conclusion}

This thesis investigated uncertainty quantification methods for Graph Neural Network surrogates of agent-based transport models. Working with 1,000 of 10,000 available MATSim scenarios (10\% subset) on a Paris road network dataset, we compared MC Dropout, ensemble variance, and conformal prediction as uncertainty quantification approaches.

\section{Summary of Findings}

\subsection{Answers to Research Questions}

\textbf{RQ1:} \emph{How effectively does MC Dropout capture epistemic uncertainty in GNN-based traffic surrogates?} MC Dropout with $S = 30$ stochastic forward passes achieves Spearman~$\rho = 0.4820$ between predicted uncertainty and absolute prediction error for Trial~8 (100 test graphs, 3,163,500 predictions). This moderate positive correlation is practically meaningful: it enables selective prediction that reduces MAE by 41.2\% at 50\% retention and error detection with AUROC of 0.76 at the top-10\% error threshold. Per-graph analysis confirms that this ranking quality is consistent across all 100 test scenarios ($\rho$ mean $= 0.464$, std $= 0.023$), with no graph below $\rho = 0.41$.

\textbf{RQ2:} \emph{How does MC Dropout compare to ensemble-based uncertainty quantification?} MC~Dropout ($\rho = 0.4908$, Experiment~A with corrected model loading; Section~\ref{sec:ensemble_caveat}) outperforms seed-based ensemble variance ($\rho = 0.4370$) by 12.3\% and multi-model ensemble disagreement ($\rho = 0.4333$, Experiment~B) by 13.3\%.  Combined MC~Dropout + ensemble variance ($\rho = 0.4909$) provides negligible improvement over MC~Dropout alone.  These results confirm that MC~Dropout is the most effective single-method uncertainty estimator in this setting, at substantially lower computational cost than maintaining multiple trained models.

\textbf{RQ3:} \emph{Does combining uncertainty estimates from architecturally identical models provide benefits?} Marginally.  Combining MC~Dropout and seed-based ensemble variance (Experiment~A, $\rho = 0.4909$) barely exceeds MC~Dropout alone ($\rho = 0.4908$), indicating that ensemble variance is largely redundant once MC~Dropout is available.  Multi-model ensemble disagreement (Experiment~B, $\rho = 0.4333$) provides useful but weaker uncertainty ranking than MC~Dropout.  The weighted ensemble prediction ($R^2 = 0.5656$) does not improve over the best individual model (T8, $R^2 = 0.5957$).  Architecturally diverse ensembles were not tested in this work.

\textbf{RQ4:} \emph{Can distribution-free methods transform raw uncertainty into trustworthy prediction intervals?} Yes. Split conformal prediction achieves near-exact marginal coverage (90.02\% at 90\% nominal, 95.01\% at 95\% nominal) on 1,581,750-node evaluation sets. Post-hoc temperature scaling ($T = 2.70$) reduces ECE from 0.265 to 0.048---an 82\% improvement---though residual miscalibration persists at high nominal levels. Adaptive conformal prediction, which normalises nonconformity scores by MC Dropout $\sigma$, narrows the conditional coverage range from [62.9\%, 98.6\%] to [90.0\%, 96.2\%] across uncertainty deciles, demonstrating that the MC Dropout ranking signal carries genuine calibration information. These three post-hoc methods---temperature scaling, standard conformal prediction, and adaptive conformal prediction---require no model retraining and form a practical calibration hierarchy.

\subsection{Primary Results}

\begin{enumerate}
    \item \textbf{Trial~8 achieves the best predictive performance among comparable trials:} $R^2 = 0.5957$, MAE $= 3.96$ veh/h, RMSE $= 7.12$ veh/h (test set, 3,163,500 predictions, 80/10/10 split). Trial~8 uses the \texttt{GATConv(64$\to$1)} final layer and dropout $= 0.2$.

    \item \textbf{MC Dropout provides effective uncertainty ranking:} Trial~8 MC Dropout ($S = 30$, 100 test graphs) achieves Spearman~$\rho = 0.4820$ between predicted uncertainty and absolute prediction error. This confirms that high uncertainty reliably predicts where the model makes larger errors.

    \item \textbf{Selective prediction substantially reduces effective error:} Filtering to the 50\% most confident predictions reduces MAE by 41.2\% (from 3.95 to $\approx$2.32~veh/h, MC Dropout mean). This demonstrates practical value: for decisions where reliability is critical, uncertainty-guided filtering provides meaningful accuracy improvements.

    \item \textbf{Conformal prediction achieves guaranteed coverage:} Split conformal prediction on T8 (50/50 calibration--evaluation split, 50 graphs each) delivers 90.02\% coverage at $\pm$9.92~veh/h and 95.01\% at $\pm$14.68~veh/h, confirming that coverage-guaranteed intervals are achievable for traffic prediction. A separate calibration audit using a larger 20/80 split corroborates these results (90.2\% at $\pm$9.99~veh/h; 95.1\% at $\pm$14.77~veh/h; 2,530,800 evaluation nodes), demonstrating robustness to the choice of calibration set size.

    \item \textbf{MC Dropout uncertainty is not a calibrated standard deviation:} The factor $k_{95} = 11.34$ required to achieve 95\% coverage via $\pm k\sigma$ intervals (vs $k = 1.96$ for a calibrated Gaussian) confirms that raw MC Dropout uncertainty should be used as a ranking signal only, not as a probability distribution.

    \item \textbf{Ensemble methods provide useful but weaker uncertainty ranking:} MC~Dropout ($\rho = 0.4908$) outperforms seed-based ensemble variance ($\rho = 0.4370$) and multi-model disagreement ($\rho = 0.4333$).  Combined MC~Dropout + ensemble variance ($\rho = 0.4909$) provides negligible additional benefit.  The weighted multi-model ensemble ($R^2 = 0.5656$) does not improve over the best individual model (T8, $R^2 = 0.5957$).  A PyTorch Geometric API version mismatch originally corrupted these results; the corrected evaluation (Section~\ref{sec:ensemble_caveat}) enables a reliable comparison.

    \item \textbf{UQ findings are robust across trials:} Replication of selective prediction and calibration audit on Trial~7 ($\rho = 0.4460$, dropout $= 0.3$) confirms the same qualitative conclusions: MC Dropout provides useful uncertainty ranking, raw $\sigma$ values are not calibrated, and conformal prediction restores near-exact coverage. The main findings are not specific to a single model configuration (Section~\ref{sec:t7_crosscheck}).

    \item \textbf{Proper scoring rules confirm sharp but underdispersed predictive distributions:} The Continuous Ranked Probability Score (CRPS/MAE $= 0.857$) indicates that the Gaussian predictive distribution is sharper than the deterministic baseline. However, Probability Integral Transform (PIT) analysis reveals systematic underdispersion (KS statistic $= 0.245$, excess mass in the first and last histogram bins), confirming that the MC Dropout $\sigma$ underestimates the true prediction error spread. Winkler score comparisons show that conformal intervals substantially outperform na\"{i}ve Gaussian intervals (32.3 vs 49.7 at the 90\% level), quantifying the benefit of distribution-free calibration.

    \item \textbf{$S = 30$ forward passes provide near-optimal uncertainty estimates:} An $S$-convergence study on 10 test graphs shows that increasing $S$ from 30 to 50 improves Spearman~$\rho$ by only 1.0\% (0.4584 to 0.4632), with the convergence curve flattening sharply after $S \approx 25$. This confirms that $S = 30$ represents a practical cost--quality trade-off for MC Dropout uncertainty estimation.
\end{enumerate}

\subsection{Practical Recommendations}

For practitioners deploying GNN-based traffic surrogates, MC Dropout with $S = 30$ samples and dropout $= 0.2$ provides the best uncertainty ranking found in this work. Uncertainty estimates should be used for selective prediction---retaining the 50\% most confident predictions reduces MAE by 41.2\%---and conformal prediction should be applied when formal coverage guarantees are contractually or operationally required. Raw MC Dropout $\sigma$ should not be interpreted as a calibrated standard deviation.  Multi-model ensembles provide complementary but weaker uncertainty ranking ($\rho = 0.4333$ vs $0.4908$ for MC~Dropout); architecturally diverse ensembles have not been tested here but could provide stronger uncertainty signals. The cross-trial replication on T7 confirms that these recommendations are not specific to a single hyperparameter configuration.

\section{Contributions}

The primary contributions of this thesis are:

\begin{enumerate}
    \item \textbf{Systematic UQ evaluation for GNN traffic surrogates:} First comprehensive comparison of MC Dropout, ensemble variance, and conformal prediction for node-level traffic volume change prediction at scale (3.16 million predictions per evaluation).

    \item \textbf{Empirical evidence for MC Dropout effectiveness:} Demonstrated $\rho = 0.4820$ uncertainty-error correlation for T8 MC Dropout. T8 (dropout $= 0.2$, 80/10/10 split) achieves higher $\rho$ than T5--T7 ($\rho$: 0.4186--0.4460, dropout $= 0.3$). Note that T8 differs from T5--T6 in both dropout rate and data split, and from T7 in dropout rate and learning rate; batch size is identical (8) across all four trials, so the $\rho$ improvement cannot be attributed solely to the reduced dropout rate.

    \item \textbf{Conformal prediction coverage verification:} Confirmed that split conformal prediction achieves nominal coverage guarantees on large-scale traffic prediction (1.58M-node evaluation sets).

    \item \textbf{Selective prediction utility:} Quantified that uncertainty-guided filtering achieves 41.2\% MAE reduction at 50\% retention, providing concrete operational guidance.

    \item \textbf{Identification of ensemble evaluation limitations:} Diagnosed a PyTorch Geometric API version mismatch that silently corrupts \texttt{GATConv} weights when loading checkpoints with \texttt{strict=False}, providing a cautionary finding for reproducibility in GNN ensemble research.

    \item \textbf{Proper scoring rule evaluation of MC Dropout predictive distributions:} Applied CRPS, PIT, and Winkler score analysis to characterise the MC Dropout predictive distribution beyond simple ranking metrics, revealing systematic underdispersion and quantifying the calibration benefit of conformal prediction over na\"{i}ve Gaussian intervals.
\end{enumerate}

\section{Broader Impact}

This work addresses the question of how to quantify prediction uncertainty in GNN-based transport surrogates. As cities increasingly rely on machine learning for policy decisions, understanding model limitations becomes necessary. Uncertainty quantification provides a mechanism for expressing model confidence, enabling more informed decision-making.

Our findings demonstrate that effective UQ is achievable with minimal architectural changes (MC Dropout, dropout $= 0.2$) and manageable computational overhead, and that formal coverage guarantees are attainable through conformal prediction. These methods require no specialised infrastructure and can be applied to existing trained models.

\section{Future Work}

This work leaves several open questions. The most direct extension is training on the \textbf{full 10,000-scenario dataset} to assess whether model accuracy and uncertainty quality scale proportionally with training data. On the architectural side, \textbf{diverse ensembles} combining GNN variants with different inductive biases (GAT, GraphSAGE, GCN) may provide the complementary uncertainty signals that identical-architecture ensembles failed to deliver here. Within the UQ methods themselves, \textbf{conformal risk control}~\cite{angelopoulos2023conformal} offers a path to tighter, adaptive coverage guarantees beyond the fixed-width intervals produced by standard split conformal prediction. \textbf{Post-hoc calibration} of MC Dropout uncertainty---through isotonic regression or temperature scaling---could reduce $k_{95}$ from 11.34 toward 1.96, allowing $\sigma$ to serve as a genuine probability estimate rather than a ranking signal. Finally, \textbf{multi-city validation} across different network topologies and traffic regimes, combined with explicit \textbf{epistemic--aleatoric decomposition}---a direction attempted via heteroscedastic NLL in this work (Trial~9, $R^2 = 0.02$; Section~\ref{sec:trial9_heteroscedastic}) but requiring stronger approaches such as evidential deep learning~\cite{amini2020deep,seitzer2022pitfalls}---would strengthen confidence in the generalisability of these findings.

\section{Closing Remarks}

Deploying machine learning surrogates responsibly in urban planning requires more than accurate predictions---it requires systematic quantification of where those predictions may fail. Uncertainty quantification provides that quantification in a form that is operationally useful: MC Dropout gives planners a per-prediction confidence score, conformal prediction gives them a coverage guarantee, and selective prediction lets them trade breadth of coverage for depth of accuracy. The experiments here, conducted on a 10\% data subset under realistic resource constraints, show that all three are achievable with the PointNetTransfGAT architecture at no meaningful cost to predictive accuracy. The practical case for uncertainty-aware surrogate deployment in transport planning is supported by the experiments presented here.
